\section{Thiết kế chương trình}

\subsection{Cấu trúc tổng quát}

Chương trình Battleship được chia thành ba giai đoạn chính:

\begin{enumerate}
    \item \textbf{Menu và chọn chế độ}: Hiển thị menu chính, cho phép người chơi chọn Play Game, View Instructions, hoặc Quit. Nếu chọn Play, người chơi chọn One Player (vs Computer) hoặc Two Players.
    
    \item \textbf{Setup Phase}: Mỗi người chơi (hoặc Computer tự động) đặt 6 tàu lên bàn cờ. Mỗi tàu phải đáp ứng các điều kiện hợp lệ. Nếu không hợp lệ, yêu cầu nhập lại.
    
    \item \textbf{Battle Phase}: Hai người chơi lần lượt tấn công. Mỗi lượt, người chơi nhập tọa độ, hệ thống kiểm tra trúng/trượt, cập nhật bàn cờ, rồi kiểm tra điều kiện thắng. Trò chơi kết thúc khi một bàn cờ không còn tàu.
\end{enumerate}

\subsection{Biểu diễn dữ liệu}

Bàn cờ 7×7 được ánh xạ thành mảng 1D với 49 phần tử:

$$\text{index} = \text{row} \times 7 + \text{col}$$

Mỗi phần tử là một 32-bit word. Byte offset để truy cập:

$$\text{byte\_offset} = \text{index} \times 4$$

Ví dụ: Ô (3, 3) có index = 3×7+3 = 24, byte\_offset = 24×4 = 96.

\subsection{Các mảng chính}

\begin{table}[H]
\centering
\small
\begin{tabular}{|l|l|l|}
\hline
\textbf{Mảng} & \textbf{Kích thước} & \textbf{Mục đích} \\
\hline
boardP1 & 49 words & Trạng thái tàu Player 1 (1=có, 0=không) \\
\hline
boardP2 & 49 words & Trạng thái tàu Player 2 (1=có, 0=không) \\
\hline
memP1 & 49 words & Lịch sử bắn Player 1 (0=chưa, 1=miss, 2=hit) \\
\hline
memP2 & 49 words & Lịch sử bắn Player 2 (0=chưa, 1=miss, 2=hit) \\
\hline
inputBuffer & 64 bytes & Lưu input từ người dùng \\
\hline
autoBoardTemplate & 49 words & Board mẫu của Computer (one player mode) \\
\hline
gameMode & 1 word & Loại chế độ chơi (1=One Player, 2=Two Players) \\
\hline
\end{tabular}
\caption{Các mảng dữ liệu chính}
\end{table}

\subsection{MIPS Implementation Notes}

Chương trình sử dụng các khái niệm MIPS sau:

\begin{itemize}
    \item \textbf{Arithmetic}: Tính index mảy từ tọa độ, tính độ dài tàu.
    \item \textbf{Memory access}: `lw` đọc giá trị ô, `sw` ghi giá trị ô.
    \item \textbf{Branches}: `beq`, `bne`, `blt` để kiểm tra điều kiện.
    \item \textbf{Procedures}: `jal` gọi thủ tục, `jr \$ra` quay về, stack lưu `\$ra` khi thủ tục lồng nhau.
    \item \textbf{Syscalls}: In/đọc chuỗi, in/đọc số, đọc ký tự.
\end{itemize}
